\section{Step 28 Finish Maxval and enter raster output}\label{ppm-s28}

Store the three ASCII digits of 255 and its LF separator. Advance the cursor to the first raster position only after those four bytes are present. Do not add an extra whitespace byte to make the raster visually separate.

Author region \texttt{e\_maxval}, keeping its BEGIN/END markers and the supplied caller, checkpoint labels and cleanup paths. Stop at \texttt{after\_e\_maxval} for the observation.

Use recipe \texttt{s28-1} with the cumulative source through this edit.

\subsection{Expected observation}
The complete worked header has eleven bytes, and the next destination byte is its first sample position.

\Needspace{9\baselineskip}
\subsection{Check your understanding}
How is Maxval 255 represented in the header?

\begin{itemize}
  \item \textbf{A} One binary byte 255.
  \item \textbf{B} Three ASCII digits.
  \item \textbf{C} A four-byte color word.
\end{itemize}
